\section*{第 \ref{ch:b2:polymer-synthesis} 章——聚合物：合成、结构与性质}

\begin{solution}{exo:b2:polymer-synthesis:1}
聚乙烯 \ce{-[CH2-CH2]-}；聚丙烯 \ce{-[CH2-CH(CH3)]-}；聚氯乙烯
\ce{-[CH2-CHCl]-}；聚苯乙烯 \ce{-[CH2-CH(C6H5)]-}；尼龙-6,6 \ce{-[NH(CH2)6NH-CO(CH2)4CO]-}。
\end{solution}

\begin{solution}{exo:b2:polymer-synthesis:2}
PET：\omterm{def:b2:polymer-synthesis:step-chain}{逐步聚合}（生成聚酯，放出水）。聚苯乙烯：\omterm{def:b2:polymer-synthesis:step-chain}{链式聚合}。尼龙-6,6：\omterm{def:b2:polymer-synthesis:step-chain}{逐步聚合}。
聚甲基丙烯酸甲酯：\omterm{def:b2:polymer-synthesis:step-chain}{链式聚合}（\omterm{def:b2:polymer-synthesis:polymer}{单体}含 \ce{C=C}）。聚乳酸由羟基酸生成：
\omterm{def:b2:polymer-synthesis:step-chain}{逐步聚合}，\omterm{def:b2:polymer-synthesis:polymer}{单体}为 \ce{A-B} 型。
\end{solution}

\begin{solution}{exo:b2:polymer-synthesis:3}
$\bar M_n = 1000 \times \qty{104.15}{g/mol} \approx \qty{104}{kg/mol}$（苯乙烯 \ce{C8H8}，
\qty{104.152}{g/mol}）。
% ledger: aw:C, aw:H
\end{solution}

\begin{solution}{exo:b2:polymer-synthesis:4}
\omterm{def:b2:polymer-synthesis:tacticity}{间同立构}。是的：规整的链能够堆砌成微晶。
\end{solution}

\begin{solution}{exo:b2:polymer-synthesis:5}
质量分数均为 0.5。$1/\bar M_n = 0.5/10 + 0.5/100 = 0.055$，$\bar M_n \approx
\qty{18.2}{kg/mol}$；$\bar M_w = 0.5 \times 10 + 0.5 \times 100 = \qty{55}{kg/mol}$；$\text{\DH} \approx
3.0$。
\end{solution}

\begin{solution}{exo:b2:polymer-synthesis:6}
$1/(1-0.98) = 50$；$1/(1-0.995) = 200$。
\end{solution}

\begin{solution}{exo:b2:polymer-synthesis:7}
$r = 1/1.01$；当 $p = 1$ 时，$\bar X_n = (1+r)/(1-r) = (1.01 + 1)/(1.01 - 1) = 201$。
\end{solution}

\begin{solution}{exo:b2:polymer-synthesis:8}
$R_p \propto [\mathrm I]^{1/2}$：变为原来的 $\sqrt2 \approx 1.41$ 倍。$\nu \propto [\mathrm I]^{-1/2}$：
除以 $\sqrt 2$，约为原值的 0.71 倍。
\end{solution}

\begin{solution}{exo:b2:polymer-synthesis:9}
纯聚氯乙烯的\omterm{def:b2:polymer-synthesis:transitions}{玻璃化转变温度}高于室温：是刚性的玻璃态。小分子增塑剂
位于链与链之间，将它们隔开，使链段在较低温度下即可运动：$T_g$ 降至
室温以下，材料在使用时呈橡胶态、柔韧。
\end{solution}

\begin{solution}{exo:b2:polymer-synthesis:10}
$\dfrac{\mathrm d}{\mathrm dx}\bigl(x p^{\,x-1}\bigr) = p^{\,x-1}(1 + x\ln p)$，当 $x = -1/\ln p$ 时为零
（为极大值：符号由 $+$ 变为 $-$）。当 $p$ 接近 1 时，$\ln p \approx -(1-p)$，故 $x \approx
1/(1-p) = \bar X_n$。
\end{solution}

\begin{solution}{exo:b2:polymer-synthesis:11}
第一阶段，\omterm{def:b2:acyl-substitution:transesterification}{酯交换}：\ce{C6H4(COOCH3)2 + 2HOCH2CH2OH -> C6H4(COOCH2CH2OH)2 + 2CH3OH}，
蒸出甲醇。第二阶段：每个二酯分子带有两个 \ce{CH2CH2OH} 端基；一端进攻另一分子的酯基，
并放出一分子乙二醇，随即抽出。
该\omterm{def:b2:polymer-synthesis:polymer}{单体}以恰当比例同时带有两种基团，与 \ce{A-B} 型\omterm{def:b2:polymer-synthesis:polymer}{单体}一样，平衡
$r = 1$ 已自动满足，移除乙二醇则使 $p$ 趋向 1。
\end{solution}

\begin{solution}{exo:b2:polymer-synthesis:12}
等摩尔投料：比值 $= (2.0 + 1)/(1 + 0.5) = 2$，即\omterm{def:b2:polymer-synthesis:polymer}{单体} 1 的单元数为\omterm{def:b2:polymer-synthesis:polymer}{单体} 2 的两倍。当 $r_1 = r_2 =
0$ 时，比值 $= [\mathrm M_1][\mathrm M_2]/([\mathrm M_2][\mathrm M_1]) = 1$，与投料比无关：每个自由基
只加成另一种\omterm{def:b2:polymer-synthesis:polymer}{单体}，形成交替\omterm{def:b2:polymer-synthesis:copolymer}{共聚物}。
\end{solution}

\begin{solution}{pb:b2:polymer-synthesis:1}
\textbf{1.} \ce{H2N(CH2)6NH2} 与 \ce{HOOC(CH2)4COOH}。
\textbf{2.} 每形成一个键放出一分子水：
\[\mathrm{{-}COOH + H_2N{-} \longrightarrow {-}CO{-}NH{-} + H_2O}.\]
\textbf{3.} 盐中每个二胺恰好对应一个二酸；若分别称量两种液体或固体，
则达不到卡罗瑟斯方程所要求的精确配比。
\textbf{4.} \ce{-[NH(CH2)6NH-CO(CH2)4CO]-}，\ce{C12H22N2O2}，\qty{226.32}{g/mol}。
\textbf{5.} \omterm{def:b2:polymer-synthesis:polymer}{重复单元}含两个由\omterm{def:b2:polymer-synthesis:polymer}{单体}衍生的结构单元：$M_0 = 226.32/2 = \qty{113.16}{g/mol}$。
\textbf{6.} 二胺中的碳原子（6 个）与二酸中的碳原子（6 个）。
\textbf{7.} $\bar X_n = 50$，$\bar M_n = 50 \times 113.16 \approx \qty{5.66}{kg/mol}$。
\textbf{8.} $\bar X_n = 100$，$\bar M_n \approx \qty{11.3}{kg/mol}$。
\textbf{9.} 当 99~\% 的基团发生反应时，平均链仅含 100 个结构单元，对于强纤维而言链太短；
之后每继续接近完全转化一步，链长便翻倍或增至三倍。
\textbf{10.} $\bar X_n = 10\,000/113.16 = 88.4$；$p = 1 - 1/88.4 \approx 0.989$。
\textbf{11.} 通过除去水实现：将熔体加热至远高于熔点，最后在减
压下进行，使平衡不断向酰胺方向移动。
\textbf{12.} 单官能度杂质会封闭链端，而任何杂质都会破坏 1\,:\,1 的配比。
\textbf{13.} $r = 1/1.01 \approx 0.990$。
\textbf{14.} $\bar X_n = (1+r)/(1-r) = 201$；$\bar M_n \approx 201 \times 113.16 \approx
\qty{22.7}{kg/mol}$。
\textbf{15.} $\bar X_n = 1.9901/(1.9901 - 2 \times 0.9901 \times 0.99) \approx 1.9901/0.0297 \approx
67$。
\textbf{16.} 胺基：当酸基耗尽时，每条链都以 \ce{NH2} 结尾。
\textbf{17.} 它将胺端基转化为不能继续反应的酰胺基，起封端作用，限制 $\bar M_n$，
并在\omterm{def:b2:polymer-synthesis:polymer}{聚合物}重新熔融时阻止进一步增长。
\textbf{18.} 为保证纺丝和模塑所需熔体黏度可重复，否则
链在熔体中继续反应，黏度会随之漂移。
\textbf{19.} $\text{\DH} = 1 + p = 1.99$。
\textbf{20.} $\bar M_w = 1.99 \times 11.3 \approx \qty{22.5}{kg/mol}$。
\textbf{21.} 数目分数 $n_1 = 1 - p = 0.01$，即 1~\% 的分子；质量分数 $w_1 = (1-p)^2 =
10^{-4}$，即 0.01~\% 的质量。
\textbf{22.} $x = -1/\ln 0.99 \approx 99.5$ 附近，即接近 $\bar X_n = 100$（见练习 10）。
\textbf{23.} $1000/226.32 = \qty{4.42}{mol}$ 的\omterm{def:b2:polymer-synthesis:polymer}{重复单元}，每个含两个酰胺键，共
\qty{8.84}{mol} 水，$8.84 \times 18.015 \approx \qty{159}{g}$。
\textbf{24.} 链太短则纤维强度弱；链太长则熔体黏度过大，难以
通过喷丝头的细孔。
\textbf{25.} $\bar X_n = 20\,000/113.16 = 176.7$，$p = 1 - 1/176.7 = \boldsymbol{\approx 0.994}$。
\end{solution}
